%% subcaptionlite.tex -- documentation for the subcaptionlite package
%%
%% Copyright (c) 2026 Munehiro Yamamoto <munepixyz@gmail.com>
%% License: LPPL 1.3c or later
%%
%% Build: pdflatex subcaptionlite.tex (twice)
\documentclass{article}
\usepackage[margin=28mm]{geometry}
\usepackage[T1]{fontenc}
\usepackage{lmodern}
\usepackage{graphicx}
\usepackage{xcolor}
\usepackage{subcaptionlite}
\usepackage[colorlinks, allcolors=blue!50!black]{hyperref}
\usepackage[capitalise]{cleveref}
\newcommand*\sclversion{v1.0 (2026/10/26)}
\newcommand*\cs[1]{\texttt{\char`\\#1}}
\newcommand*\pkg[1]{\textsf{#1}}
\newcommand*\env[1]{\texttt{#1}}
\newcommand*\meta[1]{\textnormal{$\langle$\textit{#1}$\rangle$}}

\title{The \pkg{subcaptionlite} package\\
  \large Subfigures and subcaptions without the caption package}
\author{Munehiro Yamamoto\\
  \url{https://github.com/munepi/subcaptionlite}}
\date{\sclversion}

\begin{document}
\maketitle

\section{Introduction}

The \pkg{subcaption} package needs the \pkg{caption} package, which
redefines \cs{@makecaption} of the class.  Thus it cannot be used with
a class that has its own design of captions: loading it changes the
captions of the whole document.

The \pkg{subcaptionlite} package provides the \env{subfigure}
environment, \cs{subcaption} and \cs{subcaptionbox} with the same
syntax as \pkg{subcaption}, without touching \cs{@makecaption} of the
class.  \Cref{fig:example} is made with it.

\begin{figure}[htbp]
  \centering
  \begin{subfigure}{.3\linewidth}
    \centering\includegraphics[width=.9\linewidth]{example-image-a}
    \caption{The subfigure environment}\label{fig:example-a}
  \end{subfigure}\hfill
  \subcaptionbox{\cs{subcaptionbox} with a width\label{fig:example-b}}%
    [.3\linewidth]{\includegraphics[width=.9\linewidth]{example-image-b}}\hfill
  \subcaptionbox{\cs{subcaptionbox} without a width\label{fig:example-c}}{%
    \includegraphics[width=.25\linewidth]{example-image-c}}
  \caption{Three subfigures}\label{fig:example}
\end{figure}

\paragraph{Supported formats/engines}
pdf\LaTeX, Xe\LaTeX, Lua\LaTeX{} (with or without Lua\TeX-ja),
p\LaTeX{} and up\LaTeX{} (with dvipdfmx).  The package requires
\LaTeX2e{} 2020-10-01 or later and loads \pkg{etoolbox}.

\section{Usage}

\begin{verbatim}
\usepackage{subcaptionlite}

\begin{figure}
  \centering
  \begin{subfigure}{.3\linewidth}
    \centering\includegraphics[width=.9\linewidth]{example-image-a}
    \caption{The subfigure environment}\label{fig:example-a}
  \end{subfigure}\hfill
  \subcaptionbox{\cs{subcaptionbox} with a width\label{fig:example-b}}%
    [.3\linewidth]{\includegraphics[width=.9\linewidth]{example-image-b}}\hfill
  \subcaptionbox{\cs{subcaptionbox} without a width\label{fig:example-c}}{%
    \includegraphics[width=.25\linewidth]{example-image-c}}
  \caption{Three subfigures}\label{fig:example}
\end{figure}
\end{verbatim}

\begin{description}
\item[\texttt{\char`\\begin\{subfigure\}[\meta{pos}]\{\meta{width}\}}]
  A \env{minipage} of \meta{width} (\meta{pos}: \texttt{c} by default),
  inside which \cs{caption} and \cs{caption*} mean \cs{subcaption} and
  \cs{subcaption*}.
\item[\texttt{\char`\\subcaption\{\meta{text}\}}, \texttt{\char`\\subcaption*\{\meta{text}\}}]
  A numbered or unnumbered subcaption, e.g.\ at the bottom of a
  \env{minipage}.  It is centered and set in \cs{small}.
\item[\texttt{\char`\\subcaptionbox\{\meta{caption}\}[\meta{width}]\{\meta{contents}\}}]
  Puts \meta{contents} centered in a \env{minipage} of \meta{width},
  followed by \cs{subcaption}\texttt{\{\meta{caption}\}}.  Without
  \meta{width}, the natural width of \meta{contents} is used.  If
  \meta{contents} is wider than \meta{width}, the \env{minipage} keeps
  \meta{width}.  Boxes are aligned at the bottom.
\item[\texttt{\char`\\captionsetup[\meta{type}]\{\meta{options}\}}]
  Accepted and ignored, so that the settings for the \pkg{caption}
  package can be left in the document.
\end{description}

\paragraph{Numbering and references}
The counter \texttt{subfigure} is reset by \texttt{figure}.
\cs{thesubfigure} is \texttt{(\cs{alph}\{subfigure\})} and
\cs{p@subfigure} is \cs{thefigure}, so that \cs{ref} gives
``\ref{fig:example-a}'' and \cs{Cref} of \pkg{cleveref} gives
``\Cref{fig:example-b}''.  With \pkg{hyperref}, each subfigure has its
own anchor (\cs{theHsubfigure}).

The \cs{caption} of the figure may come after the subcaptions (as in
\cref{fig:example}) or before them.  When a subcaption comes first, it
steps the figure counter in advance and the \cs{caption} that follows
does not step it again, as the \pkg{subcaption} package does; without
this, the subfigures would get the number of the previous figure.  A
figure with subcaptions and no \cs{caption} is numbered as well.

\section{Customization}

The following macros can be redefined in the class or in a local style
(\cs{makeatletter} is needed in the preamble of a document).

\begin{center}
\begin{tabular}{lp{.55\linewidth}}
  \hline
  Macro & Default\\
  \hline
  \cs{subcaption@font} & \cs{small}\\
  \cs{subcaption@sep}  & the space between the label and the text:
    \verb|\hspace{.5zw}| with p\TeX, \verb|.5\zw| with Lua\TeX-ja (or a
    class that defines \cs{zw}), \verb|.5em| otherwise\\
  \cs{subcaption@skip} & the space between the subfigure and its
    subcaption: \verb|\vskip.25\baselineskip|\\
  \hline
\end{tabular}
\end{center}

\noindent
The engine is checked when the subcaption is typeset, so Lua\TeX-ja may
be loaded after \pkg{subcaptionlite}.  The label itself is
\cs{thesubfigure}.  For example,
\begin{verbatim}
\makeatletter
\renewcommand*\subcaption@font{\footnotesize\sffamily}
\renewcommand*\subcaption@sep{\enspace}
\makeatother
\renewcommand\thesubfigure{\alph{subfigure})}
\end{verbatim}

\section{Limitations}

\begin{itemize}
\item Subcaptions are for figures only: there is no subtable.
\item \pkg{subcaptionlite} cannot be used together with the
  \pkg{caption} or \pkg{subcaption} package.
\end{itemize}

\section{Tests}

\verb|make test| compiles the test files twice with pdf\LaTeX,
Xe\LaTeX, Lua\LaTeX, p\LaTeX{} and up\LaTeX{} (+ dvipdfmx) and checks
the log.  The helpers in \texttt{test-common.tex} check the label that
\cs{label} would write, the references, the counters, the widths of
the boxes, the anchors of \pkg{hyperref}, and that \cs{@makecaption} of
the class is unchanged; the log must have \texttt{SCL-RESULT OK} and no
error or warning.

\begin{center}
\begin{tabular}{lp{.6\linewidth}}
  \hline
  Test & What it covers\\
  \hline
  \texttt{test-basic.tex}    & \env{subfigure}, \cs{subcaption},
                               \cs{subcaptionbox}, the caption of the
                               figure below/above the subfigures or none,
                               \env{figure*}, \cs{ref}, the label
                               separator\\
  \texttt{test-hyperref.tex} & \texttt{test-basic.tex} with
                               \pkg{hyperref} (loaded after and before
                               the package) and \pkg{cleveref}\\
  \texttt{test-japanese.tex} & Japanese classes (ltjsarticle, jsarticle,
                               ujarticle), Lua\TeX-ja loaded after the
                               package\\
  \hline
\end{tabular}
\end{center}

\section{History}

\begin{description}
\item[v1.0 (2026/10/26)] First public release.  Support every engine;
  fix the figure number when the caption of the figure comes before
  the subcaptions, and after a figure whose subfigures follow its
  caption.
\end{description}

\section{License}

This package is licensed under the conditions of the \LaTeX{} Project
Public License, either version~1.3c of this license or any later
version.

\medskip
\noindent
Copyright \copyright{} 2026 Munehiro Yamamoto

\end{document}
